\documentclass[10pt,a4paper]{amsart}
\usepackage[margin=19mm]{geometry}
\usepackage{amsmath,amssymb,mathtools}
\usepackage[scr=rsfs,cal=euler]{mathalfa}
\usepackage{xfrac}
\usepackage[unicode,hidelinks,bookmarks=false]{hyperref}
\newcommand{\ie}{i.e.\ }
\newcommand{\cf}{cf.\ }
\newcommand{\resp}{respectively\ }
\newcommand{\Subsection}{\subsection}
\providecommand{\nobreakdash}{\nobreak\hbox{-}}
\setlength{\emergencystretch}{2em}
\multlinegap0pt
\usepackage{parskip}
 \setlength{\parindent}{20pt}
\usepackage{thmtools}
\usepackage{cleveref}
\usepackage{enumitem}
\usepackage{color}
\newtheorem{theo}{Theorem}[section]
\newtheorem{prop}[theo]{Proposition}
\newtheorem{lemma}[theo]{Lemma}
\newtheorem{cor}[theo]{Corollary}
\newtheorem{conj}[theo]{Conjecture}
\newtheorem{defn}[theo]{Definition}
\newtheorem{rem}[theo]{Remark}
\newtheorem{eg}[theo]{Example}
\newtheorem{alg}[theo]{Algorithm}
\newcommand{\mc}[1]{\mathcal{#1}}
\newcommand{\mb}[1]{\mathbb{#1}}
\newcommand{\nib}[1]{\noindent {\bf #1}}
\newcommand{\nim}[1]{\noindent {\em #1}}
\newcommand{\brac}[1]{\left( #1 \right)}
\newcommand{\bracc}[1]{\left\{ #1 \right\}}
\newcommand{\bsize}[1]{\left| #1 \right|}
\newcommand{\brak}[1]{\left[ #1 \right]}
\newcommand{\bgen}[1]{\left\langle #1 \right\rangle}
\newcommand{\sgen}[1]{\langle #1 \rangle}
\newcommand{\bfl}[1]{\left\lfloor #1 \right\rfloor}
\newcommand{\bcl}[1]{\left\lceil #1 \right\rceil}
\newcommand{\sub}{\subseteq}
\newcommand{\subn}{\subsetneq}
\newcommand{\sups}{\supseteq}
\newcommand{\nsub}{\notsubseteq}
\newcommand{\lra}{\leftrightarrow}
\newcommand{\Lra}{\Leftrightarrow}
\newcommand{\Ra}{\Rightarrow}
\newcommand{\sm}{\setminus}
\newcommand{\ov}{\overline}
\newcommand{\ul}{\underline}
\newcommand{\wt}{\widetilde}
\newcommand{\eps}{\varepsilon}
\newcommand{\es}{\emptyset}
\newcommand{\pl}{\partial}
\newcommand{\wg}{\wedge}
\newcommand{\la}{\leftarrow}
\newcommand{\ra}{\rightarrow}
\newcommand{\da}{\downarrow}
\newcommand{\ua}{\uparrow}
\newcommand{\car}{\circlearrowright}
\newcommand{\ova}{\overrightarrow}
\newcommand{\aA}{\alpha}
\newcommand{\bB}{\beta}
\newcommand{\gG}{\gamma}
\newcommand{\dD}{\delta}
\newcommand{\iI}{\iota}
\newcommand{\kK}{\kappa}
\newcommand{\zZ}{\zeta}
\newcommand{\lL}{\lambda}
\newcommand{\tT}{\theta}
\newcommand{\sS}{\sigma}
\newcommand{\oO}{\omega}
\newcommand{\ups}{\upsilon}
\newcommand{\cE}{\mathcal E}
\newcommand{\OO}{\Omega}
\newcommand{\TT}{\Theta}
\newcommand{\Ss}{\Sigma}
\newcommand{\Ups}{\Upsilon}
\newcommand{\pk}{\textcolor{blue}}
\newcommand{\on}{\operatorname}
\begin{document}
\title[On Kahn's flow conjecture]{On Kahn's flow conjecture}
\author{Peter Keevash}
\date{}
\maketitle
\noindent\emph{Source and licence.} On Kahn's flow conjecture. Version arXiv:2609.23595v1 (CC BY 4.0). This material is distributed under the Creative Commons Attribution 4.0 International licence, http://creativecommons.org/licenses/by/4.0/, as recorded in the arXiv OAI licence metadata for this exact version and shown on the abstract page https://arxiv.org/abs/2609.23595v1. Full rights notice: licenses/arxiv-2609.23595-CC-BY-4.0.txt.\par\smallskip
\noindent\textit{Editorial scope.} Counted unit: the original \texttt{prop} environment \texttt{prop:halmos} at source lines 183--209 together with its complete proof at lines 211--283; the delivered document keeps source lines 174--284 so that every referenced label and every citation resolves inside it. Native editable source is the paper's own concatenated TeX; the AMS wrapper is editorial formatting only. No publisher class, upstream script or unrelated artwork is redistributed.
\section{Geometry of two subspaces}

We start by formulating Halmos' theory, including a proof for the reader's convenience.
Fix subspaces $U,V\sub\mb R^m$. Let $P,Q$ be orthogonal projections onto $U,V$.
For a unit vector $u\in U$, the angle $\theta$ it makes with $V$ satisfies
$\cos^2 \theta = \|Qu\|^2 = \langle u,Qu\rangle = \langle u,Cu\rangle$, where $C:=PQP|_U$. 
Thus the geometry is described by the spectral theory of $C$, as follows.



\begin{prop}\label{prop:halmos}
The $1$- and $0$-eigenspaces of $C$ are $U \cap V$ and $U \cap V^\perp$.
Let \( U_*= U\cap\bigl((U\cap V)\oplus(U\cap V^\perp)\bigr)^\perp\).
Then $U_*$ is $C$-invariant and has an orthonormal eigenbasis
$u_1,\ldots,u_r$ with \(Cu_j=\lambda_j u_j\),
where each \( \lambda_j \in (0,1) \).
Writing $c_j=\sqrt{\lambda_j}$ and $s_j=\sqrt{1-\lambda_j}$, there are
orthonormal vectors $v_1,\ldots,v_r\in U^\perp$ such that the spaces
$L_j=\operatorname{span}\{u_j,v_j\}$ are pairwise orthogonal and are
preserved by both $P$ and $Q$.  In the basis $(u_j,v_j)$,
\begin{equation}\label{eq:halmos-block}
 P=
 \begin{pmatrix}
  1&0\\
  0&0
 \end{pmatrix},
 \qquad
 Q=
 \begin{pmatrix}
  c_j^2&c_js_j\\
  c_js_j&s_j^2
 \end{pmatrix}.
\end{equation}
Thus the corresponding principal angle $\theta_j\in(0,\pi/2)$ is
determined by $c_j=\cos\theta_j$, or equivalently
$\lambda_j=\cos^2\theta_j$.
\end{prop}

\begin{proof}
We note that $C$ is self-adjoint and preserves $U$.
If $u\in U$ is a unit eigenvector with $Cu=\lambda u$ then 
\(  0 \le \lambda
 =\langle u,Cu\rangle
 =\langle u,Qu\rangle
 =\langle Qu,Qu\rangle
 =\|Qu\|^2 \le \|u\|^2 = 1 \).
If $\lambda=0$, then $Qu=0$, and hence $u\in U\cap V^\perp$;
conversely, if $u\in U\cap V^\perp$ then $Cu = PQu=0$.
If $\lambda=1$ then as $\|u\|^2=\|Qu\|^2+\|(I-Q)u\|^2$
and $\|u\|^2=\|Qu\|^2=1$ we have $(I-Q)u=0$, so $Qu=u$ and $u\in U\cap V$. 
Conversely, $C$ is the identity on $U\cap V$.  

Since $C$ is self-adjoint, it preserves $U_*$ and the
spectral theorem  gives an orthonormal eigenbasis $u_1,\ldots,u_r$ 
with eigenvalues  $\lL_1,\ldots,\lL_r$ in $(0,1)$. We define
\[
 v_j=\frac{Qu_j-\lambda_j u_j}{c_js_j}.
\]
Then
\(
 P(Qu_j-\lambda_j u_j)
 =Cu_j-\lambda_j u_j
 =0
\)
and
\(
 \|Qu_j-\lambda_j u_j\|^2
 =
 \|Qu_j\|^2
 -2\lambda_j\langle u_j,Qu_j\rangle
 +\lambda_j^2\|u_j\|^2
 =\lambda_j(1-\lambda_j),
\)
so $v_j$ is a unit vector in $U^\perp$.
Now
\(
 Qu_j=c_j^2u_j+c_js_jv_j
\)
by definition, and applying $Q$ gives
\(
 Qu_j =c_j^2Qu_j+c_js_jQv_j.
\)
Substituting and rearranging gives
\(
 c_js_jQv_j
 =c_j^2s_j^2u_j+c_js_j^3v_j,
\)
so
\(
 Qv_j=c_js_ju_j+s_j^2v_j.
\)
Since also $Pu_j=u_j$ and $Pv_j=0$, both projections preserve
$L_j$, and their matrices are those in \eqref{eq:halmos-block}.

It remains to check for $j \ne k$ that $L_j$ and $L_k$ are orthogonal. 
We already have shown orthogonality between all pairs in $\{u_j,v_j,u_k,v_k\}$
except $(v_j,v_k)$. Note that
\(
 \langle Qu_j,Qu_k\rangle
 =\langle u_j,Qu_k\rangle
 =\lambda_k\langle u_j,u_k\rangle
 =0.
\)
Substituting $Qu_j=\lambda_j u_j+c_js_jv_j$
and $Qu_k=\lambda_k u_k+c_ks_kv_k$
and expanding gives 
\(
 c_js_jc_ks_k\langle v_j,v_k\rangle=0,
\)
so $v_j\perp v_k$.
\end{proof}

\end{document}
